% The binomial scale itself is the target of the base-four congruence.
\section{A common witness for the binomial scale and exponent}
\label{sec:common-witness}

The first exponent test can use the binomial scale $U$ itself as its
target. That value is already computed, so the separate product $bw$
disappears. This change preserves the unit-scale norm and its index
modulus. Combined with the coefficient encoding of
Section~\ref{sec:borrow-mask}, it gives 102 operations.

Use the fixed coefficient index of Theorem~\ref{thm:best103}, including
its ordinary targets $4F_i^h+\delta^2$, special target $\delta^2$,
and exponent $L=2^{32}$. Retain the notation
\[
\begin{gathered}
 N=n,\quad R=r,\quad U=wN^2,\quad Y_P=sN^2,\\
 a_0=Y_P(U+1),\quad A=a_0+4,\quad
 D_P=UY_P,\quad Q=UY_P^2,\quad P=2Q+1,\quad J=2R+1.
\end{gathered}
\]
The positive supplied unknown $a$ still means $a_0$.
Keep the shared norm, first index equation and interval
\begin{equation}\label{eq:common-witness-pell}
\begin{gathered}
 \tau(\tau+1)=(D_P^2+U)(Y_Pk)^2,\qquad
 k=R+1+hD_P,\\
 c=kY_P+\eta,\qquad k=\eta+\zeta,
 \qquad\tau,h,\eta,\zeta>0.
\end{gathered}
\end{equation}
Change only E14 among the supplied-coordinate Pell polynomials,
replacing its target $bw$ by $U$:
\begin{equation}\label{eq:common-witness-exponent}
 d=U+c(A-4)+\gamma(8A-17).
\end{equation}
No parameter or witness is added. The defining identity $U=wN^2$
and the index modulus $D_P=UY_P$ are unchanged.

\subsection{Exact indices before the exponent equations}

The smaller positive offset in Section~\ref{sec:borrow-mask} retains
all preliminary coding bounds:
\[
\begin{gathered}
 N=q^8,\qquad q>4b\ge8,\qquad 3L\le B\le N\le R,\\
 0<S<N,\quad0\le T<N,\qquad
 R\le2N^3-2N^2<2N^3.
\end{gathered}
\]
These are established without using any Pell equation. Hence
$N,R\ge64$, $U,Y_P\ge N^2$, $a_0>N^4>J$, and
$D_P\ge N^4>R+1$.

The first index equation gives $k\ge R+2$, and the interval gives
$c>Y_Pk>J$. The signed block therefore gives
$c=\psi_A(J)$ and $d=\chi_A(J)$.
The norm in~\eqref{eq:common-witness-pell} is the triangular form
of the ordinary Pell equation with parameter $P=2Q+1$, so
$k=\psi_P(t)$ for some positive index $t$.
Since $P-1=2D_PY_P$ is divisible by $D_P$, the Pell congruence
and $D_P>R+1$ imply
\[
 t=R+1+vD_P,\qquad v\ge0.
\]
If $v\ge1$, then $D_P>R$ and Pell growth gives
\[
 \frac ck\le\frac{(2A)^{2R}}{(2P-1)^{R+D_P}}
 \le\left(\frac{2A}{2P-1}\right)^{2R}<\frac12.
\]
The last inequality follows from
$(2P-1)-4A=4Y_P[U(Y_P-1)-1]-15>0$.
This contradicts $c/k>Y_P$. Thus the exact indices are
\begin{equation}\label{eq:common-witness-indices}
 c=\psi_A(J),\quad d=\chi_A(J),\quad
 k=\psi_P(R+1),\quad 2\tau+1=\chi_P(R+1).
\end{equation}
The change to E14 has not entered this argument.

\subsection{Recovering the radices in the required order}

Put $\xi=(U+1)^{2R}/U^R$. The unit-scale lower estimate at
\eqref{eq:common-witness-indices} gives
\[
 \frac ck\ge
 \xi\left(1+\frac7{2a_0}\right)^{2R}
     \left(1+\frac1{2Q}\right)^{-R}>\xi,
\]
using $14Q>a_0$. The interval gives $\xi<Y_P+1$, so
\begin{equation}\label{eq:common-witness-size}
 Y_P\ge U^R,\qquad A>a_0=Y_P(U+1)>U^{R+1}.
\end{equation}
As in Section~\ref{sec:unit-scale}, the independent bound
$R<2N^3$ gives $8R/a_0<16/N<1/2$. The upper growth estimate and
$(1+t)^m\le(1-mt)^{-1}<1+2mt$ therefore yield
\begin{equation}\label{eq:common-witness-error}
 0<\frac ck-\xi<\frac{32R}{U+1}.
\end{equation}
This error need not be bounded by an absolute constant yet.

The changed exponent target satisfies
\[
 4^{3J}=64\cdot4096^R
 \le64N^{2R}\le64U^R<U^{R+1}<A,
 \qquad U^3\le U^R<A.
\]
The source $\chi$ exponent criterion applied to
\eqref{eq:common-witness-exponent} now gives
\begin{equation}\label{eq:common-witness-power}
 U=4^J.
\end{equation}
In particular $U>64R$, so~\eqref{eq:common-witness-error} is
strictly less than $1/2$.

The new first exponent conclusion determines the radices in a
different order. Since $N^2$ divides the power of two $U=wN^2$,
$N$ is a power of two. Since $N=q^8$, so is $q$.
No conclusion about $b$ is needed yet.

The unchanged second norm, gap and index congruence identify
$\kappa=\psi_A(L)$ and $\mu=\chi_A(L)$, since $0<L<J<A$.
Their exponent test has the size bounds
\[
 B^{3L}\le B^B\le N^N\le U^R<A,
 \qquad q^3\le N\le R<U^R<A.
\]
It therefore gives $q=B^L$. Now $B$ is a power of two, because
$q$ is a power of two and $L>0$. Finally $B=Hb^2$ implies that
$b$ is a power of two. Thus all binary radix facts are known before
the first-mask decoding uses them.

The binomial expansion has integer part $F=\lfloor\xi\rfloor$ with
\[
 0<\xi-F<\frac{2^{2R}}U<\frac14,
 \qquad F\equiv\binom{2R}{R}\pmod U.
\]
The interval and the ratio bounds give
\[
 F<\xi<\frac ck<Y_P+1,
 \qquad Y_P<\frac ck<\xi+\frac12<F+\frac34.
\]
Integrality forces $Y_P=F$. The unchanged definitions
$U=wN^2$ and $Y_P=sN^2$ now imply
\begin{equation}\label{eq:common-witness-binomial}
 N^2\mid\binom{2R}{R}.
\end{equation}

\subsection{Composition with the new coefficient encoding}

The conclusions just obtained are exactly those needed by
Section~\ref{sec:borrow-mask}: $b,q,N$ are powers of two,
$q=B^L$, and~\eqref{eq:common-witness-binomial} supplies binary
nonoverlap of the packed blocks. Positivity of the new offset then
gives
\[
 \Omega>Z\lambda/2,\qquad
 \mathcal C^2<2Bq/Z<q^2,
 \qquad\mathcal C=x+g.
\]
The first two masks recover the possible quotient ambiguity
$\ell=\ell_0+mq$, $e=e_0-m$ and the bounds required by
Lemma~\ref{lem:borrow-high-window}. That lemma forces $m=0$.
The low coefficients have incoming carries in $\{-1,0\}$; the
special row and guard force $\delta=1$, and every ordinary row
then forces $F_i^h=0$. These decoding arguments do not use E14
internally. They apply unchanged once its replacement has supplied
the required exponent and binomial conclusions, proving sufficiency.

For necessity, first use the new canonical coding construction of
Section~\ref{sec:borrow-mask}. It has $\delta=1$, $u=x^2$, zero
dummy coordinates and a sufficiently large power of two $b$.
Its changed $\sigma$ is positive because
\[
 S_3=\lambda(Bq-Z\mathcal C^2)+2e\mathcal C^2>0.
\]
The ordinary and special target digits are one or two, so the masks
pass. Form its resulting $S,T,R$. This supplies the packed bounds
and~\eqref{eq:common-witness-binomial} before choosing Pell witnesses.

Now set
\[
\begin{gathered}
 J=2R+1,\quad U=4^J,\quad w=U/N^2,\quad
 Y_P=\lfloor\xi\rfloor,\quad s=Y_P/N^2,\\
 a_0=Y_P(U+1),\quad A=a_0+4,\quad Q=UY_P^2,
 \quad P=2Q+1,\\
 c=\psi_A(J),\quad d=\chi_A(J),\quad k=\psi_P(R+1).
\end{gathered}
\]
The quotient $w$ is positive integral: canonical $N$ is a power of
two and $N^2\le R^2\le4^R<4^J=U$. The binomial expansion and
divisibility make $s$ positive integral. The same estimates at these
indices give $Y_P<c/k<Y_P+3/4$, so choose the positive integers
\[
 \eta=c-kY_P,\quad \zeta=k-\eta,\quad
 \tau=\frac{\chi_P(R+1)-1}{2},\quad
 h=\frac{k-R-1}{UY_P}.
\]
Odd $P$ gives the positive integrality of $\tau$. The congruence
modulo $P-1=2UY_P^2$ and strict Pell growth give that of $h$.
This construction needs no parity condition on $R$.

Choose the remaining Pell witnesses at the new $A$ by the generic
constructions already established in Section~\ref{sec:main-base-four}.
After choosing positive $f,i$ for E16, put
\[
 G=1+(A+1)(f^2-1),\qquad
 o=\frac{\chi_G(J)+d}{f},\qquad
 j=\frac{\psi_G(J)-J}{c},
\]
and take $\kappa=\psi_A(L)$, $\mu=\chi_A(L)$,
$\Delta=(\kappa-L)/(A-1)$ and $\phi=c-\kappa$.
Their congruences and strict growth make every required witness
positive integral. The exponent quotients are positive because
\[
 d-(A-4)c>3c>U,\qquad
 \mu-(A-B)\kappa>(B-1)\kappa>q,
\]
using $A>U^3,q^3$ and $J,L\ge2$; their moduli are positive.
Thus the new E14 and every other equation hold. The coding witnesses
are chosen first, and all dependent Pell coordinates are then chosen
for their new packed $R$.

\begin{theorem}\label{thm:best102}
For each r.e.\ set, the fixed coefficient encoding of
Section~\ref{sec:borrow-mask} gives a membership-equivalent system
in 34 positive unknowns and 22 equations whose straight-line
certificate has \textbf{102 arithmetic instructions}:
56 multiplications and 46 additions. Fixed numerals, supplied
parameters and equality tests are free.
\end{theorem}
\begin{proof}
The preceding arguments establish the complete positive-domain
equivalence and the order of the two constructions.
Starting from the 104-operation schedule, the offset change deletes
the addition $q+1$ and uses $(B\lambda)q$ in its place.
The common-witness change deletes the product $bw$ and uses the
existing register $U=wN^2$ directly in E14. These edits concern
disjoint register definitions, so they remove one addition and one
multiplication without introducing an instruction, unknown or equation.

The resulting count is 56 multiplications and 46 additions.
The rebuilt coefficient index affects only supplied constants.
The checker \file{round22\_1980\_composed\_certificate.py}
verifies the complete primitive list, all source residuals and both
triangular corrections. It also checks that the two schedule edits
commute and that both discarded registers are absent.
\end{proof}
